\section{Formal Definitions}
\label{app:math}

Let $G=(V,E)$ be the interactome with $N=|V|$ nodes, and $d(u,v)$ the shortest-path length (number of edges) between $u$ and $v$.

\paragraph{Closest distance.} For a drug with target set $T$ and disease module
$S$~\cite{guney2016network}:
\begin{equation}
  d_c(S,T) = \frac{1}{|T|}\sum_{t\in T}\min_{s\in S} d(s,t).
\end{equation}

\paragraph{Degree-matched randomization.} Nodes are grouped into bins of similar degree, each bin holding at least $b_{\min}$ nodes (e.g., 100). A random sample replaces every node of $T$ and of $S$ by a uniformly drawn node from the same bin. Over $R$ samples, with mean $\mu$ and standard deviation
$\sigma$ of $d_c$,
\begin{equation}
  z = \frac{d_c(S,T)-\mu}{\sigma},
\end{equation}
with an empirical one-sided $p$-value $p = (1+\#\{r: d_c^{(r)}\le d_c\})/(1+R)$.

\paragraph{Separation.} For two node sets $A$ and
$B$~\cite{menche2015uncovering},
\begin{equation}
  s_{AB} = \langle d_{AB}\rangle - \frac{\langle d_{AA}\rangle + \langle d_{BB}\rangle}{2},
\end{equation}
where $\langle d_{AA}\rangle$ is the mean over $a\in A$ of the distance to the nearest \emph{other} member of $A$, and $\langle d_{AB}\rangle$ is the mean, over all nodes of $A$ and $B$, of the distance to the nearest node of the other set. $s_{AB}<0$ means the two target sets overlap in the network.

\paragraph{Complementary exposure.} A drug pair is favored when both drugs are close to the disease ($z<0$) and their target sets are separated ($s_{AB}\ge0$)~\cite{cheng2019network}.

\paragraph{Combination QUBO (optional).} For a shortlist of $n$ drugs with $x_i=1$ if drug $i$ is chosen, let $w_i=-z_i>0$ reward disease proximity and $P_{ij}=\max(0,-s_{ij})$ penalize overlapping target sets. Choosing $k$ drugs:
\begin{equation}
  \min_{x\in\{0,1\}^n}\; -\sum_i w_i x_i + \lambda\sum_{i<j}P_{ij}x_ix_j
  + A\Big(\sum_i x_i - k\Big)^2 .
\end{equation}
Expanding with $x_i^2=x_i$ gives $x^{\top}Qx$ (upper-triangular form) plus the constant $Ak^2$, with
\begin{equation}
  Q_{ii} = -w_i + A(1-2k), \qquad Q_{ij} = \lambda P_{ij} + 2A \quad (i<j).
\end{equation}
$A$ must be large enough that violating the cardinality constraint never pays off (e.g., $A > \max_i w_i + \lambda\sum_j P_{ij}$); $\lambda$ is swept to show the trade-off between proximity and overlap.

%======================================================================
\section{Kernel Design Details}
\label{app:kernels}

\subsection{Data layout}

\begin{itemize}[nosep]
  \item Graph: CSR arrays \texttt{rowptr} ($N+1$ ints) and \texttt{col} ($2E$ ints).
  \item Distance matrix: $D$ as \texttt{uint8\_t}, row-major, $N^2$ bytes (324~MB at $N=18{,}000$); 255 marks ``unreachable''.
  \item MS-BFS state per batch: \texttt{seen}, \texttt{frontier}, \texttt{next}, each $N\times W$ 64-bit words ($W=16$ gives 1{,}024 sources per batch; 3 arrays $\times$ 18{,}000 $\times$ 16 $\times$ 8~B $\approx$ 7~MB).
  \item Target sets and the disease module: flattened index arrays with offsets; degree bins: node lists per bin, plus a node-to-bin map.
\end{itemize}

\subsection{Bit-parallel multi-source BFS}

\begin{algorithm}[h]
\caption{One batch of MS-BFS (sources $s_0,\dots,s_{64W-1}$)}
\begin{algorithmic}[1]
\State $\text{seen}[s_j] \gets \text{frontier}[s_j] \gets$ bit $j$ set; $D[s_j][s_j]\gets 0$
\For{$\ell = 1, 2, \dots$}
  \ForAll{nodes $v$ and words $w$ \textbf{in parallel}} \Comment{one thread per $(v,w)$}
    \State $m \gets 0$
    \ForAll{$u \in \text{neighbors}(v)$} \State $m \gets m \;|\; \text{frontier}[u][w]$ \EndFor
    \State $\text{next}[v][w] \gets m \;\&\; \lnot\,\text{seen}[v][w]$
    \ForAll{set bits $j$ of $\text{next}[v][w]$} \State $D[s_{64w+j}][v] \gets \ell$ \EndFor
  \EndFor
  \State $\text{seen} \gets \text{seen} \,|\, \text{next}$;\; swap(frontier, next)
  \If{no bit was set this level} \textbf{break} \EndIf
\EndFor
\end{algorithmic}
\end{algorithm}

This is a ``pull'' formulation: each thread reads its neighbors' words and
writes only its own, so no atomics are needed. Threads for consecutive $v$
and the same $w$ write consecutive bytes of row $s_{64w+j}$ of $D$, which
coalesces. Iterating set bits uses \texttt{\_\_ffsll} and clearing the lowest
bit. High-degree nodes cause load imbalance; if profiling shows this, nodes
are sorted by degree and heavy nodes are assigned a warp each. The number of
levels equals the largest distance reached in the batch (typically
$\le 15$).

\subsection{Proximity kernel}

One block per drug; thread $t$ processes samples $r = t, t+B, t+2B,\dots$ For
each sample it seeds a Philox counter-based generator with (drug ID, sample
ID), so the CPU reference reproduces the exact same random nodes. It draws
the random target and disease sets from the degree bins, evaluates $d_c$ with
reads from $D$, and accumulates $\sum d_c$ and $\sum d_c^2$ in registers. A
CUB block reduction gives $\mu$ and $\sigma$. Reads from $D$ are scattered;
the real disease module's $|S|$ rows (e.g., 100 $\times$ 18{,}000 bytes
$\approx$ 1.8~MB) fit in L2 cache, which helps the non-random term.

\subsection{Exact QUBO solver}

Thread $t$ owns the $2^m$ patterns whose top $n-m$ bits equal $t$. It
evaluates the first pattern in full ($O(n^2)$), then walks the rest in
Gray-code order: flipping bit $i$ changes the cost by
$\Delta = (1-2x_i)\big(Q_{ii} + \sum_{j\ne i} c_{ij}x_j\big)$, where $c_{ij}$
is the coefficient of $x_ix_j$ (i.e., $Q_{ij}$ or $Q_{ji}$, whichever is in
the upper triangle), which is $O(n)$, or $O(1)$ if a running vector of local fields
is kept. Warp and block minimum reductions find the best pattern. At $n=28$
this is $2.7\times10^8$ patterns.

\subsection{Work estimates}

\begin{center}
\small
\begin{tabular}{@{}lll@{}}
\toprule
Step & Sequential cost & Example size \\
\midrule
All-pairs BFS & $O(N(N+E))$ & $\approx 4.5\times10^{9}$ edge visits (Hetionet); $\approx 2.4\times10^{10}$ (BioGRID) \\
Proximity, all drugs & $O(n_{\text{drug}}\,R\,|T|\,|S|)$ & $\approx 2.6\times10^{8}\times\overline{|T|}$ look-ups (1{,}383 drugs, $R$=1k, $|S|$=185) \\
Separation, shortlist & $O(n^2\,|T|^2)$ & small \\
QUBO brute force & $O(2^n n)$ & $\approx 7.5\times10^{9}$ at $n$=28 \\
\bottomrule
\end{tabular}
\end{center}

%======================================================================
\section{Optional Quantum Part: QAOA with cuQuantum}
\label{app:quantum}

\emph{Note: the extension is now framed as a solver benchmark in which QAOA is one
solver among several; see Appendix~\ref{app:extension}. The material below remains the
technical description of the QAOA part.}

Each shortlisted drug is one qubit. The QUBO is mapped to an Ising cost
Hamiltonian $H_C$ by $x_i=(1-Z_i)/2$. QAOA of depth $p$ prepares
\begin{equation}
  |\psi(\gamma,\beta)\rangle = \prod_{\ell=1}^{p} e^{-i\beta_\ell H_M}
  e^{-i\gamma_\ell H_C}\,|+\rangle^{\otimes n}, \qquad H_M=\sum_i X_i,
\end{equation}
and a classical optimizer (COBYLA) tunes the $2p$ angles to minimize
$\langle\psi|H_C|\psi\rangle$~\cite{farhi2014qaoa}. The circuit is built in
CUDA-Q and simulated with cuStateVec~\cite{bayraktar2023cuquantum}; the
expectation is computed exactly from the state vector, and the final state is
sampled to read out combinations. A state of $n$ qubits needs $2^n$ complex
numbers: 2~GiB at $n=28$ (complex64), 8~GiB at $n=30$.

Reported for $p=1,\dots,5$ and $n\in\{12,16,20,24,28\}$:
\begin{itemize}[nosep]
  \item approximation ratio $r = (E_{\text{worst}} - E_{\text{QAOA}})/(E_{\text{worst}} - E_{\text{opt}})$
        over valid solutions, with $E_{\text{opt}}$ from the GPU brute force;
  \item probability of sampling the optimum;
  \item fraction of samples violating the ``exactly $k$ drugs'' constraint;
  \item comparison with simulated annealing~\cite{kirkpatrick1983optimization}
        given the same number of cost evaluations.
\end{itemize}
No quantum speedup is claimed: at these sizes brute force is exact and fast.
The goal is to measure QAOA solution quality on realistic, biology-derived
instances.

%======================================================================
\section{Evaluation Protocol}
\label{app:eval}

\subsection{Timing}

\begin{itemize}[nosep]
  \item Sequential baseline: single-threaded C++ implementing the identical
        algorithm (standard queue-based BFS from each node; proximity loop
        with the same RNG). An OpenMP version is reported as an extra
        baseline.
  \item Scaling: network size $N$ varied through Hetionet, BioGRID, and (if time) STRING confidence cut-offs
        (e.g., 900, 700, 400); random samples
        $R\in\{100,1{,}000,10{,}000\}$.
  \item Protocol: two warm-up runs, then the median of 10. GPU time by CUDA
        events, reported both kernel-only and end-to-end including
        transfers. Speedup $=T_{\text{CPU}}/T_{\text{GPU}}$ per step and
        overall. Nsight Compute reports achieved memory bandwidth, cache hit
        rates, occupancy, and warp divergence.
\end{itemize}

\subsection{Correctness}

\begin{itemize}[nosep]
  \item $D_{\text{GPU}} = D_{\text{CPU}}$ exactly (byte-wise comparison);
        NetworkX as an independent check on small graphs.
  \item With shared Philox seeds, $z$-scores from GPU and CPU agree to
        floating-point rounding.
  \item Brute-force QUBO optimum checked against exhaustive CPU search for
        small $n$.
\end{itemize}

\subsection{Biological validation}

\begin{itemize}[nosep]
  \item \textbf{Positives:} drugs that have entered AD clinical
        trials~\cite{cummings2024pipeline}. \textbf{Metrics:} area under the
        ROC and precision--recall curves, and enrichment of positives in the
        top 50/100 (hypergeometric test).
  \item \textbf{Baselines:} random ranking; ranking by number of targets;
        ranking by target degree; ranking by direct overlap between targets
        and AD genes.
  \item \textbf{Circularity check:} many AD trial drugs target amyloid or tau
        directly, which are in the AD module. Metrics are also reported after
        removing drugs whose targets include AD-module genes, so the result
        is not just ``drugs aimed at AD genes are close to AD genes.''
  \item \textbf{Robustness (enabled by speed):} Spearman correlation of
        $z$-scores between Hetionet and BioGRID, across $R$, and across
        alternative AD modules (e.g., GWAS-only vs.\ GWAS plus known genes).
\end{itemize}

%======================================================================
\section{Limitations and Risks}
\label{app:limits}

\begin{itemize}[nosep]
  \item \textbf{Incomplete network:} well-studied proteins have more known
        interactions; degree-matched sampling reduces but does not remove
        this bias.
  \item \textbf{Label quality:} being in a clinical trial does not mean a
        drug works; positives are ``considered promising,'' not ``effective.''
  \item \textbf{Memory-bound kernels:} proximity look-ups are scattered, so
        speedup depends on cache behavior; this is measured, not assumed.
  \item \textbf{Amdahl's law:} data preparation in Python may dominate
        end-to-end time; it is reported separately.
  \item \textbf{Quantum part:} simulation only, $n\lesssim30$; QAOA may
        return invalid or suboptimal combinations at low depth, which is
        reported as a result.
  \item \textbf{Not medical advice:} candidates are hypotheses for further
        study.
\end{itemize}

%======================================================================
% Appendices F--I: decision log, team guide, revised extension.
% Written for teammates who are new to bioinformatics and quantum computing.
%======================================================================
\section{Team Guide and Decision Log}
\label{app:team}

\subsection{How to read this appendix}

Appendices A--E are the technical reference. Appendices F--I record
\emph{decisions}: what we chose, why, and what it changes. You do not need
biology or quantum background to follow them; each term is explained in
Section~\ref{sec:glossary}.

Every decision has an ID (D1, D2, \dots) and a status:
\textbf{Decided} (settled), \textbf{Proposed} (suggested, waiting for the team
to confirm), or \textbf{Open} (not answered yet). If a decision contradicts
earlier text, the ``Affects'' column says which part it replaces, and
\emph{the newer decision wins}.

\subsection{Priorities}
\label{sec:priorities}

We have two goals, in this order:
\begin{enumerate}[nosep]
  \item \textbf{A good class paper and demo} for the course (Milestones 1--3).
  \item \textbf{A publishable paper} built on the same work.
\end{enumerate}
In practice this means:
\begin{itemize}[nosep]
  \item Modules 1--4 and the CPU-vs-GPU comparison are the deliverable. They
        are finished, tested, and written up before any extra work starts.
  \item The extension (Module 5) starts only after the core results can be
        reproduced from a clean checkout. The fallback in
        Section~\ref{sec:milestones} stays: if it is not ready, it is future
        work.
  \item After the extension, hand-written AVX-512/assembly is an optional
        extra (D9). The custom state-vector kernel comes after that (D5).
  \item Work that only matters for publication (extra statistics, verifying
        citations and data) is tracked in Section~\ref{sec:pubgate} and may
        not delay a class milestone.
  \item When one choice helps both goals, take it. Fair baselines, for
        example, improve the grade \emph{and} survive peer review.
\end{itemize}

\subsection{Decision log}

\begin{table}[!ht]
\centering
\small
\begin{tabularx}{\linewidth}{@{}l>{\raggedright\arraybackslash}X>{\raggedright\arraybackslash}Xll@{}}
\toprule
ID & Decision & Why & Status & Affects \\
\midrule
D1 & Class paper first, publication second. & Course deadlines are fixed; publication needs extra rigor that can come after. & Decided & Schedule (Tab.~\ref{tab:schedule}); all optional items \\
D2 & Modules talk to each other \emph{only} through files with fixed formats (Section~\ref{sec:modules}). & Each person can build and test one module without understanding the others. & Decided & Tab.~\ref{tab:modules}; Sec.~\ref{sec:method} \\
D3 & Schedule time to teach every member the minimum they need (Section~\ref{sec:teaching}). & There are three of us, none with a bioinformatics or quantum background; a module nobody understands cannot be debugged or defended. & Decided & Schedule \\
D4 & Reframe Module 5 as a \textbf{benchmark of QUBO solvers}: exact, classical, and QAOA simulated on a GPU (Appendix~\ref{app:extension}). No quantum hardware is used. & The honest question is solution quality, not speed. & Decided & Abstract; Sec.~1.4; RQ4; App.~\ref{app:quantum} \\
D5 & The custom GPU state-vector kernel is the \textbf{last} option, started only if everything else is done. & It is the most work and the least needed for the course or the benchmark. & Decided & App.~\ref{app:quantum}; Week~12 \\
D6 & Use a ladder of CPU baselines, not one (Section~\ref{sec:ladder}). & Separates the gain from bit tricks, from SIMD, from multicore, and from the GPU. & Proposed & Sec.~\ref{sec:method}; App.~\ref{app:eval}; RQ1 \\
D7 & \textbf{Hetionet is the primary network}; BioGRID is the comparison network and the heavier timing workload (Section~\ref{sec:datasets}). & Hetionet bundles interactions, drug targets, and AD genes in one download. We stay open to switching the primary network if results or licensing demand it. & Decided & Sec.~\ref{sec:method} (Inputs); Milestone~2; App.~\ref{app:eval} \\
D8 & The disease module is randomized in every random sample, as in Appendix~\ref{app:math}. & Matches the formal definition. Whether one random module is shared across drugs for the same sample is still to be fixed (Section~\ref{sec:ablation}). & Decided & App.~\ref{app:math}; App.~\ref{app:kernels} \\
D9 & Build \textbf{CUDA C++ with a C++ CPU baseline first}, including the AVX2 baseline rung. Hand-written AVX-512/assembly is an optional target after the quantum benchmark, unless the team finds it more novel. & One toolchain for the GPU pipeline and the GPU-simulated quantum part; AVX2 is already in the baseline ladder. & Decided & Platform table; Sec.~\ref{sec:milestones} \\
\bottomrule
\end{tabularx}
\caption{Decision log.}
\label{tab:decisions}
\end{table}

\subsection{Modules, owners, and the file contract}
\label{sec:modules}

\begin{plain}
Think of the project as an assembly line. Each station (module) takes a file in,
does one job, and puts a file out. If your output file has the agreed format
and the tests pass, nobody downstream needs to know how you made it.
\end{plain}

\begin{table}[!ht]
\centering
\small
\begin{tabularx}{\linewidth}{@{}l>{\raggedright\arraybackslash}p{2.0cm}>{\raggedright\arraybackslash}X>{\raggedright\arraybackslash}X>{\raggedright\arraybackslash}p{3.1cm}@{}}
\toprule
Module & Owner & Reads & Writes & Must understand \\
\midrule
1.\ Data prep & \todo{name} & Raw downloads: interactions, drug--target pairs, AD genes & Graph (CSR) and lists of node IDs & File formats, ID mapping; \emph{no} biology beyond ``a gene has an ID'' \\
2.\ Distances & \todo{name} & Graph (CSR) & Distance matrix ($N\times N$ bytes, binary file) & BFS, bit operations, CUDA; \emph{no} biology \\
3.\ Proximity & \todo{name} & Distance matrix, drug targets, AD module & Table with $d$, $\mu$, $\sigma$, $z$, $p$ per drug & Random sampling, $z$-score \\
4.\ Ranking and validation & \todo{name} & Proximity table, list of AD-trial drugs & Ranked list and validation report & ROC/AUROC, baselines, permutation tests \\
5.\ Combinations & \todo{name} & Ranking, distance matrix & QUBO matrix, solver results & QUBO, optimization, solver comparison \\
\bottomrule
\end{tabularx}
\caption{Module ownership. There are three of us and five modules, so some
people own more than one.}
\label{tab:owners}
\end{table}

Cross-cutting roles that also need an owner: the \textbf{timing harness}
(Appendix~\ref{app:eval}), the \textbf{correctness tests}, and the
\textbf{paper}.

\paragraph{Correctness across versions.} The sequential C++ and CUDA versions of
Modules~2 and~3 read and write the \emph{same files} and are compared
byte for byte (distances) and to floating-point rounding ($z$-scores, with the
shared counter-based RNG). NetworkX provides an independent check on small
graphs (Appendix~\ref{app:eval}). A module is ``done'' when its tests pass
and its output matches the sequential version.

\subsection{Teaching plan}
\label{sec:teaching}

\begin{plain}
Ground rules: (1) always start from a toy example you can draw on paper (10
proteins, 3 drugs); (2) teach each person only what their module needs;
(3) after each session, every member explains the idea back in five minutes
to someone else. If you cannot explain it simply, we have not taught it
properly yet.
\end{plain}

\begin{table}[!ht]
\centering
\small
\begin{tabularx}{\linewidth}{@{}c>{\raggedright\arraybackslash}p{3.0cm}>{\raggedright\arraybackslash}X>{\raggedright\arraybackslash}X@{}}
\toprule
\# & Topic & Everyone leaves able to\dots & Quick check \\
\midrule
1 & The big picture: proteins as a map, drugs as stops & Explain what a drug's proximity to AD means, using the city-map picture (Section~\ref{sec:method}). & Say in two sentences why we compare against random ``fake drugs.'' \\
2 & Graphs, BFS, and the bit-parallel trick & Run BFS by hand on the toy graph, then see how one 64-bit word holds 64 searches at once. & Fill in the first two levels of the distance matrix for 3 sources. \\
3 & Random samples and the $z$-score & Compute $d$, $\mu$, $\sigma$, and $z$ for one toy drug with a handful of random samples. & Predict what happens to $\sigma$ when the number of samples grows. \\
4 & Mapping it to SIMT/SIMD & Explain ``one thread per (node, word)'' and why writes to the distance matrix coalesce. & Sketch which thread writes which byte. \\
5 & QUBO and solvers (no quantum needed) & Write a 4-drug QUBO by hand; solve it by trying all 16 choices; explain greedy and annealing. & Say what ``simulated QAOA'' does and does not claim (Section~\ref{sec:glossary}). \\
\bottomrule
\end{tabularx}
\caption{Teaching sessions (about 30--45 minutes each). Suggested timing:
sessions 1--3 before Week~7, session~4 before Week~8, session~5 before
Week~12. \todo{confirm who teaches which session}}
\label{tab:teaching}
\end{table}

Analogies (cities, delivery trucks) are for intuition only. The definitions
in Appendix~\ref{app:math} are the authority.

\subsection{Glossary}
\label{sec:glossary}

\begin{table}[!ht]
\centering
\small
\begin{tabularx}{\linewidth}{@{}l>{\raggedright\arraybackslash}X@{}}
\toprule
Term & Meaning in this project \\
\midrule
Interactome & The graph of human proteins (nodes) and known physical interactions (edges). \\
Drug target & A protein a drug binds to. Each drug has a small set of them. \\
Disease module & The set of proteins linked to Alzheimer's disease (AD). \\
Proximity $d$ & Average distance from a drug's targets to the nearest AD protein. Smaller is closer. \\
$z$-score & How unusual $d$ is compared with random look-alike drugs. Strongly negative means unusually close. \\
Drug repurposing & Testing an approved drug for a new disease. \\
QUBO & ``Pick a yes/no value for each item to minimize a quadratic score.'' Here: which drugs to choose together. \\
QAOA & A quantum algorithm for QUBOs. We run it only as a \emph{simulation} on a GPU, because no quantum computer is used. That limits us to small problems and means we claim no quantum speedup. \\
SIMT / SIMD & Many threads (GPU) or wide vector instructions (CPU) doing the same operation on different data. \\
\bottomrule
\end{tabularx}
\caption{Glossary for teammates.}
\end{table}

%======================================================================
\section{Revised Extension: Benchmarking QUBO Solvers}
\label{app:extension}

This appendix replaces the framing of Appendix~\ref{app:quantum} (decision D4).
The QAOA material there is still valid; it becomes one solver among several.
All quantum computation here is \emph{simulated on a GPU}; no quantum hardware
(QPU) is involved.

\begin{plain}
We are \emph{not} asking ``can a quantum computer pick drug combinations
better?'' We are asking: for a realistic selection problem built from
biology, how good is a simulated quantum algorithm compared with exact and
classical solvers when each gets the same amount of effort? A clear answer of
``the simple method is as good'' is a valid result.
\end{plain}

\subsection{Research questions (replace RQ4)}

\begin{enumerate}[label=\textbf{RQ4\alph*.},leftmargin=*]
  \item With the same budget (number of cost evaluations, and separately
        wall-clock time), how close does each solver come to the exact
        optimum on biology-derived QUBOs?
  \item Are biology-derived instances easier or harder than random QUBOs of
        the same size and density?
  \item Does the way we enforce ``exactly $k$ drugs'' (a penalty term vs.\
        moves that always keep exactly $k$) change which solver wins?
\end{enumerate}

\subsection{Solver ladder}

\begin{table}[!ht]
\centering
\small
\begin{tabularx}{\linewidth}{@{}c>{\raggedright\arraybackslash}p{3.4cm}>{\raggedright\arraybackslash}p{2.4cm}>{\raggedright\arraybackslash}Xl@{}}
\toprule
\# & Solver & Kind & Role in the benchmark & Status \\
\midrule
1 & Exact brute force (GPU, Gray code) & Exact & Ground truth up to $n\approx 28$ & Decided \\
2 & Random sampling & Baseline & The floor: what luck alone gives & Proposed \\
3 & Greedy & Classical & If it is near-optimal, the problem is easy; report that & Proposed \\
4 & Simulated annealing~\cite{kirkpatrick1983optimization} & Classical & Standard heuristic & Decided \\
5 & Tabu search & Classical & Optional second heuristic & Proposed \\
6 & QAOA, $p\le 3$, simulated on GPU & Simulated quantum & The solver of interest; compared against rows 1--5 & Decided \\
\bottomrule
\end{tabularx}
\caption{Solvers in the benchmark.}
\end{table}

\subsection{Experiment design}

\begin{itemize}[nosep]
  \item \textbf{Instances.} Shortlist sizes $n\in\{12,16,20,24\}$ (exact
        solutions are available well beyond this), with $k$ fixed by the team.
        Use at least 20 instances per $n$, built from different shortlists,
        seeds, or AD-module variants, so we report medians and confidence
        intervals, not single runs.
  \item \textbf{Control instances.} Random QUBOs with the same $n$, density,
        and weight scale. This answers RQ4b.
  \item \textbf{Difficulty.} Sweep $\lambda$ (Appendix~\ref{app:math}); it
        controls how much the drugs ``conflict.'' Report the optimality gap
        and how many near-optimal solutions exist.
  \item \textbf{Fairness.} Same budget for every solver. Count cost
        evaluations for algorithmic comparisons; report wall-clock time
        separately, on the same machine.
  \item \textbf{Constraint handling (RQ4c).} Penalty version for every solver.
        For greedy, annealing, and tabu search, also a version whose moves
        keep exactly $k$ drugs. For QAOA, an XY-mixer version is an optional
        extra.
  \item \textbf{Metrics.} Optimality gap of the best valid solution;
        probability of reaching the optimum; fraction of valid outputs; best
        value found vs.\ evaluations used.
\end{itemize}

\subsection{Reporting checks}

\begin{itemize}[nosep]
  \item \textbf{Always print the luck level next to a probability.} For $n$
        drugs and $k$ chosen, a uniformly random bit string is valid with
        probability $\binom{n}{k}/2^{n}$ and optimal with probability
        $1/2^{n}$. A QAOA result means little until it is compared with these.
  \item \textbf{Choose the approximation-ratio metric on purpose.} A ratio
        computed from the \emph{best sampled} solution cannot change with
        depth $p$ once the best sample is found. Prefer the expected ratio over
        the whole sampled distribution, plus the probability of the optimum.
  \item QAOA at $p\le 3$ is a weak solver on its own. Its value in the paper
        comes from the comparison with rows 1--5 of the ladder.
\end{itemize}

\subsection{Wording rules for the paper}

\begin{itemize}[nosep]
  \item Say ``simulated QAOA benchmark.'' Never ``quantum drug discovery.''
  \item State that nothing ran on quantum hardware, that results at
        $n\le 28$ say nothing about quantum hardware, that the simulation is
        noiseless, and that brute force is exact and fast at these sizes.
  \item Report negative results (for example, ``greedy matches annealing'')
        as findings, not failures.
  \item Do not claim a speedup unless we measured one against the best
        classical baseline.
\end{itemize}

\subsection{Custom state-vector kernel (last option, D5)}

Start this only if the core, the benchmark, and the write-up are complete.
The idea: precompute the cost $C[x]$ for all $2^n$ states (the brute-force
solver already produces it); the phase step is then one elementwise
multiplication $\psi[x]\leftarrow e^{-i\gamma C[x]}\psi[x]$, and the mixer is
$n$ passes over the array. Validate against cuStateVec/CUDA-Q at small $n$
before comparing speed. If it is not done, it goes in future work and nothing
else is affected.

%======================================================================
\section{CPU Baseline Ladder and Evaluation Updates}
\label{app:baselines}

\subsection{The baseline ladder (D6)}
\label{sec:ladder}

\begin{plain}
If the GPU is $100\times$ faster than a naive CPU loop but bit tricks alone
give $30\times$, then the GPU's own contribution is the rest. (Illustrative
numbers.) The ladder lets us say exactly where each speedup comes from.
\end{plain}

\begin{table}[!ht]
\centering
\small
\begin{tabularx}{\linewidth}{@{}c>{\raggedright\arraybackslash}p{4.3cm}>{\raggedright\arraybackslash}X@{}}
\toprule
Rung & Implementation & What it isolates \\
\midrule
A & C++ queue-based BFS, 1 thread & The course's ``sequential'' baseline \\
B & C++ multi-source BFS on 64-bit words, 1 thread & The effect of bit-parallelism alone \\
C & B with AVX2 (and AVX-512 if available), 1 thread & The CPU SIMD gain \\
D & C with OpenMP over nodes & A fair multicore CPU \\
E & CUDA multi-source BFS & The GPU. \textbf{The headline comparison is E vs.\ D.} \\
\bottomrule
\end{tabularx}
\caption{CPU baselines for the all-pairs distance step. The proximity step
uses the same idea: scalar, OpenMP, and optionally AVX2 gathers. Hand-written
AVX-512/assembly beyond rung C is an optional extra (D9).}
\end{table}

Fairness rules:
\begin{itemize}[nosep]
  \item Same batch size $W$ and the same node relabeling (BFS or degree
        order) on CPU and GPU.
  \item Same RNG and seeds, so all rungs compute identical results.
  \item Report kernel-only \emph{and} end-to-end time, two warm-up runs, then
        the median of 10 (Appendix~\ref{app:eval}).
  \item AVX2 tip: with $W=16$ each node holds four 256-bit registers, so the
        neighbor loop is a bitwise OR over four registers. The scalar
        step that writes distances for newly reached sources will likely be
        the CPU bottleneck, so profile it.
  \item \textbf{Size matters.} Hetionet needs about $4.5\times10^{9}$ edge
        visits for all-pairs BFS, which a good CPU does in seconds. BioGRID
        needs about $2.4\times10^{10}$. Run the timing experiments on both, and
        with $R$ up to 10{,}000, so the comparison is not decided on the
        smallest case.
\end{itemize}

\subsection{Proximity ablation: shared random module (optional)}
\label{sec:ablation}

By D8 the disease module is redrawn for every random sample. There are two
ways to organize this. In the kernel of Appendix~\ref{app:kernels}, every
(drug, sample) pair draws its own random module, costing about $|T|\cdot|S|$
lookups. Alternatively, each \emph{sample index} draws one random module shared
by all drugs. Then one multi-source BFS per sample gives every node's distance
to the nearest module protein, and each drug costs only $|T|$ lookups. At our
sizes this has about the same cost as the naive method (a small constant
factor at best), and the benefit grows with the number of drugs. It is a design
choice and an \textbf{ablation to measure on both CPU and GPU}, not an assumed
win. It also slightly changes the statistics (samples are correlated across
drugs), so the team must pick one version for the main results.

\subsection{Validation additions}

\begin{itemize}[nosep]
  \item Report how many trial drugs survive mapping onto the network.
  \item \textbf{The validation list must be verified} before results are
        reported (Section~\ref{sec:datasets}). Until then, every validation
        number is provisional.
  \item With a validation list of about 30 drugs, the AUROC of a random
        ranking has a standard deviation of roughly 0.05, so differences of a
        few hundredths are noise. Report confidence intervals and a
        permutation test; average the random baseline over many shuffles
        instead of one.
  \item Keep the circularity check: report results with and without drugs
        whose targets are themselves AD-module genes.
  \item Include the other baselines from Appendix~\ref{app:eval} (number of
        targets, target degree, direct overlap). The number-of-targets
        baseline can be strong, because drugs with many targets are more
        likely to have been tested in trials. Beating it is a real test.
        Not beating it is a valid result to report honestly (RQ3).
\end{itemize}

%======================================================================
\section{Data, Open Questions, and Publication Checklist}
\label{app:status}

\subsection{Datasets at a glance}
\label{sec:datasets}

\begin{table}[!ht]
\centering
\small
\begin{tabularx}{\linewidth}{@{}lrrrrr@{}}
\toprule
Network & Proteins & Interactions & Drugs (${\ge}1$ target) & AD genes & All-pairs BFS work ($N\cdot 2E$) \\
\midrule
Hetionet~\cite{himmelstein2017hetionet} & 15{,}139 & 147{,}149 & 1{,}383 & 185 & $\approx 4.5\times10^{9}$ \\
BioGRID~\cite{oughtred2021biogrid} & 17{,}602 & 679{,}846 & 1{,}382 & 191 & $\approx 2.4\times10^{10}$ \\
\bottomrule
\end{tabularx}
\caption{Processed datasets. Hetionet supplies interactions, DrugBank-derived
drug--target links, and AD genes in one release; the largest observed distance
is 9, so one byte per distance suffices.}
\end{table}

\begin{itemize}[nosep]
  \item \textbf{Licenses.} Hetionet combines sources with mixed licenses, and
        some (such as DrugBank) are non-commercial. We publish download and
        processing scripts, not raw or derived copies of restricted data.
        \todo{check each source's license before the repository goes public}
  \item \textbf{Validation list.} About 30 AD-trial drugs compiled from
        general knowledge and not yet verified. Each one needs a
        ClinicalTrials.gov entry (record the NCT ID) or a listing in the annual
        pipeline review~\cite{cummings2024pipeline}.
  \item \textbf{Obvious hits.} Some drugs will rank high simply because their
        targets \emph{are} AD genes. That is why the circularity check in
        Appendix~\ref{app:eval} is mandatory.
\end{itemize}

\subsection{Open questions for the team}

\begin{table}[!ht]
\centering
\small
\begin{tabularx}{\linewidth}{@{}c>{\raggedright\arraybackslash}X>{\raggedright\arraybackslash}X@{}}
\toprule
\# & Question & Why it matters \\
\midrule
Q1 & Who verifies the validation list, and by when? & Every biological result depends on it. \\
Q2 & Do we also use the genetics-based AD module as a robustness variant? & Tests whether the ranking depends on how AD is defined. \\
Q3 & Shared random module across drugs, or one per (drug, sample) (Section~\ref{sec:ablation})? & Changes the statistics and the cost. \\
Q4 & Which GPU and CPU will be used for all timings? & The platform table still has \todo{model} fields. \\
Q5 & Which $n$ and $k$ for the solver benchmark (Appendix~\ref{app:extension})? & Fixes the scope of Module~5. \\
Q6 & Is there a target venue and date for the publication goal? & Decides how much of Section~\ref{sec:pubgate} to start. \\
\bottomrule
\end{tabularx}
\end{table}

\subsection{Publication checklist (second priority)}
\label{sec:pubgate}

\begin{itemize}[nosep]
  \item Fill in the TODO references in \texttt{refs.bib}
        (\cite{csbj2023proximity,lqt2025proximity,qmi2026combinations}) from
        the original papers, and read each one before citing it as prior work.
        Confirm the newly added entries too.
  \item Confirm or soften the novelty claim in Section~\ref{sec:rrl}:
        describe how we searched, and use ``to our knowledge.'' In particular,
        check the literature on concurrent multi-source BFS on GPUs.
        \todo{verify, for example iBFS (SIGMOD 2016)}
  \item Verify the validation list (Section~\ref{sec:datasets}) and check data
        licenses before publishing anything derived from them.
  \item Meet the statistical items in Appendix~\ref{app:baselines}
        (confidence intervals, permutation tests, many random shuffles).
  \item Freeze data versions, seeds, and software versions in the repository
        so another group can reproduce the tables.
  \item Keep the ``hypotheses, not medical advice'' framing
        (Appendix~\ref{app:limits}).
\end{itemize}
